% ============================================================
%  Section 2 — Formulation statistique du problème
% ============================================================
%
%  Contenu attendu (voir sujet PDF) :
%
%  2.1 Intérêt de la formulation statistique
%      - Expliquer pourquoi reformuler en termes de paramètres (p, µ)
%        plutôt que répondre "vrai/faux" à une question intuitive
%      - Lien avec la méthode scientifique (§ 3.3 du cours)
%
%  2.2 Paramètres d'intérêt
%      - Identifier les paramètres testés : proportions π ou moyennes m
%      - Exemple : "la proportion d'étudiant·es dormant < 6h avant un DE"
%      - Exemple : "la moyenne de notes de bac en sciences"
%
%  2.3 Formulation des hypothèses statistiques
%      - Écrire H0 et H1 pour chaque test envisagé (≥ 3 tests distincts)
%      - Préciser si bilatéral ou unilatéral (Déf. 3.2 du cours)
%      - Indiquer le type de test associé :
%          · conformité de proportion → critère z (§ 3.4.1)
%          · conformité de moyenne   → critère t Student (§ 3.4.1)
%          · χ² de conformité        → critère χ² (§ 3.4.2)
%          · χ² d'indépendance       → critère χ² (§ 3.4.2)
%
% ============================================================

\section{Formulation statistique du problème}

\subsection{Intérêt de la formulation statistique}

% TODO : rédiger ici

\subsection{Paramètres d'intérêt}

% TODO : rédiger ici

\subsection{Hypothèses statistiques}

% TODO : tableau récapitulatif H0 / H1 / type de test / bilatéral ou unilatéral
